% nuc_nnforce: the nucleon-nucleon potential energy (schematic) and the Coulomb tail
% (Vol III Ch 1, fig:nuc-nnforce).  Same schematic shape as the source (figs.py):
%   V_N(r) = -60 exp(-(r-0.9)^2/0.35) MeV, plus a core 400 (0.7-r)^1.6 MeV for r < 0.7 fm.
% This is an illustration of the features (core, well, short range), not a fitted
% potential.  Coulomb energy of two protons: e^2/(4 pi eps0 r) = 1.44/r MeV (r in fm).
% Right panel: V_N passes -1.5 MeV at r = 0.9 + sqrt(0.35 ln 40) = 2.04 fm and is
% below 0.01 MeV in magnitude beyond about 2.4 fm, so at 3 fm and beyond only the
% Coulomb term is left.
\begin{tikzpicture}[declare function={vn(\r)=-60*exp(-(\r-0.9)^2/0.35);}]
\begin{groupplot}[bookplot,
    group style={group size=2 by 1, horizontal sep=19mm},
    width=6.9cm, height=5.4cm]
  \nextgroupplot[xmin=0, xmax=4.1, ymin=-75, ymax=95,
      axis x line=middle, xlabel={$r$ (fm)},
      xlabel style={at={(axis description cs:1,0.425)},anchor=north east},
      ylabel={potential energy (MeV)},
      xtick={1,2,3}, ytick={-60,-30,30,60},
      title style={font=\small}, title={two nucleons (schematic)}]
    \addplot[cPlum,domain=0.30:4,samples=160]
      {vn(x)+ (x<0.7)*400*(max(0.7-x,0.0001))^1.6};
    \node[cPlum,anchor=west,align=left] at (axis cs:0.75,62) {repulsive core};
    \node[cPlum,anchor=west,align=left] at (axis cs:1.45,-45) {strong attraction,\\ range $\approx2$\,fm};
  \nextgroupplot[xmin=0, xmax=10.4, ymin=-1.5, ymax=1.65,
      axis x line=middle, xlabel={$r$ (fm)},
      xlabel style={at={(axis description cs:1,0.46)},anchor=north east},
      ylabel={energy (MeV)},
      xtick={2,4,6,8}, ytick={-1,1},
      title style={font=\small}, title={zoomed out: only Coulomb survives}]
    \addplot[cPlum,domain=2.04:10,samples=120] {vn(x)};
    \addplot[cPrim,domain=1:10,samples=120] {1.44/x};
    \node[cPlum,anchor=west] at (axis cs:2.2,-1.2) {nuclear (any pair)};
    \node[cPrim,anchor=south west] at (axis cs:3.0,0.55) {Coulomb, $1.44/r$ (p--p only)};
\end{groupplot}
\end{tikzpicture}
